\justifying
\NSFCBodyText

\subsubsection{与本项目相关的研究工作积累}

申请人颜丰博士现为新疆大学副教授、硕士研究生导师，长期开展时间序列分析、多源信息融合、多模态视觉理解以及机器学习与深度学习研究，并具有科研技术与软件研发经历。现有积累主要来自算法与软件研究，尚未形成跨境物流现场成果；其方法可迁移到本项目三项任务：时序预测和状态空间建模用于不规则观测、动态风险与漂移分析；视觉问答、跨模态检索和小目标检测用于包装图像、安防视频与传感证据融合；软件研发及应用项目经验用于数据管线、弱网回放和边缘原型实现。

在时序建模方面，申请人以通讯作者发表 DWT-Former（\emph{Energy}, 2025），研究小波多尺度特征与 Transformer 时间表征的融合；以共同第一作者发表 WaveKAN（\emph{Neurocomputing}, 2026），研究自适应小波分解与长时序预测。

申请人主持的 XJEDU2026P005 项目聚焦多源数据驱动的光伏功率预测，创建的 ALPHA Group 也持续开展长时序预测与多源时序融合。这些工作为本项目三时钟观测、事件状态估计和仓储缓变风险建模提供方法储备。

在多模态与视觉方面，申请人在视觉问答、遥感视觉理解、跨模态检索方向形成 SPCA-Net、Knowledge-Aware Image Understanding、RSSR、RSHR+ 和 Feature Fusion Mamba Hashing 等成果，并以共同第一作者发表 S-YOLO（\emph{Neural Networks}, 2025）。申请人主持自治区天池英才“多模态融合技术研究及其在新疆产业集群中的应用”项目、自治区优秀博士后“面向通用领域的多模态视觉问答关键技术研究”，以及贴片二极管表面缺陷检测横向项目。上述事实能够证明多模态表示、视觉异常分析和模型研发能力，但不据此声称已经完成货物破损或仓储预警的前期实验。

\subsubsection{已取得的研究工作成绩}

申请人已在 \emph{Energy}、\emph{Neurocomputing}、\emph{Neural Networks}、\emph{Machine Learning}、\emph{The Visual Computer}、\emph{Computer Vision and Image Understanding}、\emph{Expert Systems with Applications} 和 \emph{IEEE Signal Processing Letters} 等期刊发表或列示相关成果，研究主题覆盖多尺度时序、长时预测、多模态交互、状态空间推理、跨模态检索和视觉检测。申请人主持项目 7 项，内容涉及多源预测、多模态融合、视觉问答、缺陷检测、数据采集和软件底座；参与低空安全管控、无人平台语义表征、公路沙害智能预警及场所安全服务等项目。上述成果构成算法研究与软件验证的证据链。2026 年论文的卷期、DOI 和最终检索状态仍以正式申报前核验结果为准。

\subsubsection{项目组能力与任务分工}

申请人具备项目总体组织、时间序列与多模态模型设计、实验评估和软件原型研发能力。项目组参与者均为研究生，按研究任务安排如下：孔哲、张建宏参与运输阶段振动、倾斜、电子封志、光感和包装图像的时序对齐、异常样本整理与模型复现；杨铄、孟星宇参与 2G/3G 弱网仿真、离线缓存与断点续传原型、集合条件证据价值策略和端侧资源评估；桑乐吕、中原参与温湿度、视频、烟雾和门禁数据的仓储风险建模、校准与预警实验；于得武参与跨任务数据规范、公共基线、分组测试与复现实验。上述内容是本项目拟执行的任务分工，不据此声称成员已具备通信硬件或物流现场专长；申请人负责总体方案、关键模型、质量控制和三项任务之间的证据状态衔接。

\subsubsection{前期基础与本项目的衔接}

前期基础与项目形成三层衔接：长时序、多尺度与状态空间建模可迁移到事件—采样—到达错位下的风险表示；多模态交互、跨模态检索和小目标检测可迁移到振动/封志/图像及温湿度/视频/门禁/烟雾的证据融合；软件技术和应用研发经历可支撑离线缓存、断点续传、价值调度与端边协同的原型实现。申请人的优势集中在算法与软件，新型通信协议和低功耗硬件不作为已有基础或原创主张，网络、能量、缓存与设备均作为可测约束处理。

申请人确认目前尚无与本申请直接对应的物流数据、代码、原型、软件著作权、专利或预实验结果。本节据此只用已公开论文、项目经历和平台条件证明相邻方法能力；直接证据将从基础可控场景、弱网模拟和可复现实验开始建立，跨领域迁移不作为已经获得的物流实验结论。

\subsubsection{研究风险及应对措施}

\textbf{技术风险：集合条件价值或风险记忆缺少独立增益。} 早期信号是预算匹配后价值策略不优于逐包强基线，或真实运输摘要与置换摘要表现相近。应对措施是保留成对重放和负对照，若假设被否证则将成果收敛为三时钟风险界与安全截止期调度；仓储主体风险校准仍独立交付，不用增加模型复杂度掩盖阴性结果。

\textbf{进度风险：真实数据、三时钟日志或现场准入晚于计划。} 早期信号是第一年度数据规范确定后仍无法获得完整采样/到达日志或授权场景。应对措施是先以可控异常和参数可追溯网络回放完成机制验证，同时持续准备有限真实场景；无货物级配对时不启动跨阶段记忆强验证，并明确降低外推论断。

\textbf{资源风险：端侧设备、功耗仪或算力服务不能覆盖预定实验。} 早期信号是无法形成可复现实测资源记录，或 A100 服务可用期/配额不能覆盖关键训练。应对措施是采用 1--2 类可获得代表性板卡、记录时延/内存/字节并把能耗明确标为估算；通过模型压缩、分阶段训练和共享计算资源完成主要假设检验，不宣称未经实测的续航倍数。
